\documentclass[a4paper, serif, 10pt]{moderncv}

%% moderncv
% style and color
\moderncvstyle{banking}
\moderncvcolor{black}

%% page layout/margins
\usepackage[top=2.5cm, bottom=2.5cm, left=1.5cm, right=1.5cm]{geometry}

\nopagenumbers{}

%% Babel config for Norwegian language
% ref:
% - https://latex3.github.io/babel/guides/locale-norwegian.html
\usepackage[norsk]{babel}

%% fontspec
\usepackage{fontspec}

\IfFontExistsTF{Linux Libertine}{
  \setmainfont[Ligatures={TeX, Common}, Numbers=OldStyle]{Linux Libertine}
}{}
\IfFontExistsTF{Linux Libertine O}{
  \setmainfont[Ligatures={TeX, Common}, Numbers=OldStyle]{Linux Libertine O}
}{}

%% CTeX
% CJK support, used to show CJK characters in the resume
%
% - fontset=none: disable builtin fontset but instead set the CJK font manually
% - heading=false: disable ctex heading
% - punct=kaiming: use kaiming punctuations styles for CJK
% - scheme=plain: use plain scheme, do not override \normalsize font size
% - space=auto: space settings for CJK characters
%
% ref:
% - http://ctan.mirrorcatalogs.com/language/chinese/ctex/ctex.pdf
\usepackage[UTF8, heading=false, punct=kaiming, scheme=plain, space=auto]{ctex}

\IfFontExistsTF{Noto Serif CJK SC}{
  \setCJKmainfont{Noto Serif CJK SC}
}{}
\IfFontExistsTF{Noto Sans CJK SC}{
  \setCJKsansfont{Noto Sans CJK SC}
}{}

%% line spacing
\usepackage{setspace}
\setstretch{1.125}

%% URL styling - use normal text instead of monospace
\urlstyle{same}

% Auto-underline all links
% Defer until \begin{document} so hyperref is already loaded
\AtBeginDocument{%
  \let\oldhref\href
  \renewcommand{\href}[2]{\oldhref{#1}{\underline{#2}}}%
}

%% Patch \httplink and \httpslink to handle full URLs
% The original moderncv macros always prepend http:// or https:// to URLs.
% This causes issues when the URL already has a protocol (e.g., https://example.com)
% resulting in malformed URLs like https://https://example.com.
% This patch checks if the URL already contains "://" and uses it directly if so.
% Note: We use \str_set:Nx (x-expansion) to expand macros like \@homepage first.
\ExplSyntaxOn
\str_new:N \g_yamlresume_url_str

\RenewDocumentCommand{\httpslink}{O{}m}{
  \str_set:Nx \g_yamlresume_url_str {#2}
  \str_if_in:NnTF \g_yamlresume_url_str {://}
    { \url{#2} }
    { \url{https://#2} }
}

\RenewDocumentCommand{\httplink}{O{}m}{
  \str_set:Nx \g_yamlresume_url_str {#2}
  \str_if_in:NnTF \g_yamlresume_url_str {://}
    { \url{#2} }
    { \url{http://#2} }
}
\ExplSyntaxOff

\name{Ingrid Solberg}{}
\title{Dataforsker med fokus på maskinlæring og prediktiv analyse}
\phone[mobile]{+47 912 34 567}
\email{ingrid.solberg@example.no}
\homepage{https://ingridsolberg.no}

\address{Karl Johans gate 1, Oslo, Oslo, Norge, 0150}{}{}

\extrainfo{{\small \faLinkedin}\ \href{https://www.linkedin.com/in/ingridsolberg}{@ingridsolberg} {} {} {} • {} {} {}
{\small \faGithub}\ \href{https://github.com/ingridsolberg}{@ingridsolberg}}

\begin{document}

\maketitle

\section{Grunnleggende informasjon}

\cvline{}{Dataforsker med master i informatikk og sterk bakgrunn innen maskinlæring,
statistikk og programvareutvikling. Har arbeidet med prediktive modeller,
kundeanalyse og stordata i finans- og telekombransjen.
%
Erfaring med hele datavitenskapens verdikjede – fra problemforståelse og
dataforberedelse til modellering, evaluering og drift i produksjon. Opptatt
av å formidle komplekse resultater på en tydelig måte og skape verdi gjennom
datadrevne beslutninger.}
\section{Utdanning}

\cventry{aug. 2017–juni 2019}
        {Master, Informatikk og data science, Poeng: A (ECTS)}
        {Universitetet i Oslo}
        {\url{https://www.uio.no/}}
        {}
        {\begin{itemize}
\item Fordypet meg i maskinlæring, statistisk læring og dyp læring med praktiske prosjekter innen tekst- og tidsserieanalyse.
\item Skrev masteroppgave om prediktive modeller for energiforbruk i bygg med gradient boosting og random forest.
\item Fikk erfaring med vitenskapelig metode, eksperimentdesign og formidling av komplekse resultater.
\item Relevant verktøykasse: Python, scikit-learn, TensorFlow, SQL og Git.
\end{itemize}
\textbf{Kurs}: Maskinlæring, Statistisk læring, Dyp læring, Datavarehus og Big Data, Tekstanalyse og NLP, Bayesiansk statistikk}

\cventry{aug. 2014–juni 2017}
        {Bachelor, Matematikk og informatikk, Poeng: A/B (ECTS)}
        {Universitetet i Bergen}
        {\url{https://www.uib.no/}}
        {}
        {\begin{itemize}
\item Ga et solid grunnlag i matematikk, statistikk og programmering med vekt på algoritmer og datastrukturer.
\item Gjennomførte bacheloroppgave i samarbeid med en lokal bedrift om visualisering av offentlige data.
\item Utviklet ferdigheter i Python, R og databasemodellering.
\end{itemize}
\textbf{Kurs}: Matematisk analyse, Lineær algebra, Algoritmer og datastrukturer, Statistikk, Databaser, Objektorientert programmering}
\section{Arbeidserfaring}

\cventry{mars 2022–Nå}
        {Dataforsker}
        {DNB}
        {\url{https://www.dnb.no/}}
        {}
        {\begin{itemize}
\item Utvikler og drifter prediktive modeller for kredittrisiko og kundeadferd i samarbeid med forretningsområder.
\item Bygger MLOps-pipeline i Google Cloud, fra datahenting til automatisk retrening og modellovervåkning.
\item Gjennomfører A/B-tester og kausale analyser for å evaluere effekten av nye produkter og kampanjer.
\item Presenterer resultater for ledelsen og bidrar til strategiske beslutninger om risiko og kundeopplevelse.
\end{itemize}
\textbf{Nøkkelord}: Prediktiv analyse, Risikomodellering, MLOps, Python, Google Cloud, A/B-testing}

\cventry{sep. 2019–feb. 2022}
        {Data Scientist}
        {Telenor}
        {\url{https://www.telenor.no/}}
        {}
        {\begin{itemize}
\item Utviklet churn-modeller og kundesegmenteringsverktøy for å redusere kundeavgang i mobil- og bredbåndsmarkedet.
\item Bearbeidet store datamengder med Apache Spark og SQL for å bygge features og analysere kundedata.
\item Samarbeidet med produkt- og markedsavdelinger for å integrere modellresultater i kampanjestrategier.
\item Etablerte rutiner for modellovervåkning og dokumentasjon, noe som forbedret teamets leveranseevne.
\end{itemize}
\textbf{Nøkkelord}: Churn-modeller, Kundesegmentering, Apache Spark, SQL, Power BI, A/B-testing}
\section{Språk}

\cvline{Norsk}{Morsmål eller tospråklig nivå \hfill \textbf{Nøkkelord}: Morsmål, Bokmål}
\cvline{Engelsk}{Fullt profesjonelt nivå \hfill \textbf{Nøkkelord}: IELTS 8.0, Akademisk skriving}
\section{Ferdigheter}

\cvline{Maskinlæring}{Ekspert \hfill \textbf{Nøkkelord}: Python, TensorFlow, PyTorch, scikit-learn, XGBoost}
\cvline{Dataanalyse og statistikk}{Ekspert \hfill \textbf{Nøkkelord}: Statistisk testing, Regresjon, Tidsserieanalyse, Bayesiansk statistikk, Pandas}
\cvline{Big Data}{Avansert \hfill \textbf{Nøkkelord}: Apache Spark, Databricks, SQL, Kafka, Airflow}
\cvline{Datavisualisering}{Avansert \hfill \textbf{Nøkkelord}: Matplotlib, Plotly, Power BI, Tableau, Seaborn}
\cvline{Programmering}{Ekspert \hfill \textbf{Nøkkelord}: Python, R, SQL, Git, FastAPI}
\section{Utmerkelser}

\cventry{nov. 2021}
        {Telenor}
        {Telenor Data Science Award}
        {}
        {}
        {Fikk pris for arbeidet med prediktive churn-modeller som bidro til bedre kundebevaring og økt lønnsomhet.}
\section{Sertifikater}

\cventry{apr. 2023}
        {Google Cloud}
        {Professional Machine Learning Engineer}
        {\url{https://cloud.google.com/learn/certification/machine-learning-engineer}}
        {}
        {}

\cventry{mai 2020}
        {Coursera / DeepLearning.AI}
        {Deep Learning Specialization}
        {\url{https://www.coursera.org/specializations/deep-learning}}
        {}
        {}
\section{Publikasjoner}

\cventry{juni 2019}
        {Prediktive modeller for energiforbruk i norske bygg}
        {Masteroppgave, Universitetet i Oslo}
        {\url{https://www.duo.uio.no/}}
        {}
        {Undersøkte hvordan gradient boosting og random forest kan forutsi timebasert energiforbruk og bidra til energieffektivisering.}
\section{Prosjekter}

\cventry{jan. 2023–juni 2023}
        {Åpen kildekode-prosjekt for prediksjon av energiforbruk basert på værdata og historiske målinger.
}
        {Energiprediksjon}
        {\url{https://github.com/ingridsolberg/energiprediksjon}}
        {}
        {\begin{itemize}
\item Utviklet en Python-pakke for tidsserieanalyse og prediksjon av strømforbruk.
\item Inneholder moduler for datahenting, feature engineering og modelltrening.
\item Dokumentert med eksempler og satt opp som FastAPI-tjeneste med Docker.
\end{itemize}
\textbf{Nøkkelord}: Python, FastAPI, XGBoost, Docker, Open source}
\section{Interesser}

\cvline{Friluftsliv}{Fjellturer, Langrenn, Kajakk}
\cvline{Teknologi}{Kunstig intelligens, Etisk AI, Datavitenskap}
\section{Frivillig arbeid}

\cventry{aug. 2021–Nå}
        {Frivillig arrangementansvarlig}
        {Oslo Data Science Meetup}
        {\url{https://www.meetup.com/Oslo-Data-Science-Meetup/}}
        {}
        {Bidrar til å arrangere faglige møter og workshops for datascience-miljøet i Oslo.
Holder introduksjonsforedrag om maskinlæring og etisk AI for nye medlemmer.}

\end{document}